\hypertarget{appendix-6ac071f8b50c}{}\label{referencestart:appendix-6ac071f8b50c}
\section{Hessians and local curvature}\label{appendix-6ac071f8b50c-hessians-and-local-curvature}

\textbf{Curriculum depth:} Developed · \textbf{First seen:} Phase 9 · \textbf{Latest development:} Phase 9

\subsection{Definition and mechanism}\label{appendix-6ac071f8b50c-definition-and-mechanism}

The Hessian collects second partial derivatives of a scalar objective. It describes how the gradient changes and supplies the quadratic term in a local Taylor model. For a smooth scalar field, the Hessian is symmetric under the usual equality-of-mixed-partials assumptions. Its eigenvectors indicate local curvature directions: positive eigenvalues give upward curvature, negative values downward curvature, and near-zero values weakly constrained directions. A positive-definite Hessian supports a strict local minimum in the unconstrained setting, while an indefinite Hessian signals a saddle-like local model. The Hessian is not the same as the curvature of an embedded geometric surface, though both involve second-order change. In optimization, exact Newton steps solve with the Hessian; practical methods may regularize or approximate it because conditioning, cost, or indefiniteness can make direct steps unreliable.

\subsection{Prerequisites}\label{appendix-6ac071f8b50c-prerequisites}

\begin{itemize}
\tightlist
\item
  \hyperref[appendix-d8a7ddb675eb]{Jacobians, pushforwards, and pullbacks}
\item
  \hyperref[appendix-50d2ea3c9b26]{Symmetric positive-definite matrices and Cholesky}
\end{itemize}

\subsection{Used by}\label{appendix-6ac071f8b50c-used-by}

\begin{itemize}
\tightlist
\item
  \hyperref[appendix-54c8f5828387]{Newton and quasi-Newton optimization}
\item
  \hyperref[appendix-2c2df8345991]{Nonlinear least squares: Gauss-Newton and Levenberg-Marquardt}
\item
  \hyperref[appendix-32f39e9a3738]{Taylor models and approximation error}
\end{itemize}

\subsection{Full teaching explanation}\label{appendix-6ac071f8b50c-full-teaching-explanation}

\begin{itemize}
\tightlist
\item
  \href{../02-Mathematics-for-Visual-Computing/main.pdf\#nameddest=02-foundations-ii-ch021-jacobians-curvature-and-local-approximation}{Phase 9 teaching section (Mathematics for Visual Computing, Mathematics for Visual Computing, Chapter 5)}
\item
  \hyperref[appendix-adf48f5dc5c1]{Visual Recall/09 - Calculus and Optimization - Visual Recall.md} · visual recall
\end{itemize}

\subsection{Source coverage}\label{appendix-6ac071f8b50c-source-coverage}

Supplied lesson parts: 826 (T04 · M050), 827 (T04 · M050).
